El procesamiento integrará los eventos por identificador de partido y orden temporal con la información de plantillas, alineaciones y posesiones. Se intentará delimitar intervalos sin sustituciones, reconstruir los cinco integrantes de cada equipo y contabilizar tiempo, marcador y posesiones durante el enfrentamiento. Los intervalos cuya reconstrucción no sea inequívoca se marcarán como no verificables.

Posteriormente se derivarán estadísticas contextuales de jugadoras que coinciden en cancha y agregados colectivos por quinteto y rival. La coincidencia contextual no se confundirá con la identificación de la defensora directa. Las características históricas se calcularán con ventanas temporales anteriores a cada partido para evitar fuga de información.

Las variables derivadas se explican en las Tablas~\ref{tab:stints_wnba} y \ref{tab:coincidencias_wnba}. El conjunto procesado servirá como entrada para la comparación de técnicas de agrupamiento, estimación del rendimiento y optimización, descritas en el Capítulo 6.